\documentclass[11pt,a4paper]{article}

\usepackage[T1]{fontenc}
\usepackage[utf8]{inputenc}
\usepackage[french]{babel}
\usepackage[margin=2.5cm]{geometry}
\usepackage{booktabs}
\usepackage{siunitx}
\usepackage{xcolor}
\usepackage{hyperref}
\usepackage{pgfplots}
\pgfplotsset{compat=1.18}
\usepackage{caption}
\usepackage{enumitem}
\usepackage{amsmath}

\sisetup{detect-weight=true, detect-family=true}

\hypersetup{
  colorlinks=true,
  linkcolor=blue!50!black,
  urlcolor=blue!50!black,
  citecolor=blue!50!black,
  pdftitle={Benchmark VASP - parallelisation NPAR/KPAR},
  pdfauthor={Pr. John Akapulco}
}

\title{\vspace{-1.5cm}Benchmark de performance VASP 6.5.0\\[2pt]
\large Réglage de la parallélisation (NPAR/KPAR) et scaling en nombre de tâches MPI}
\author{Pr. John Akapulco \\ \small \href{mailto:computchem86@gmail.com}{computchem86@gmail.com}}
\date{22 août 2026}

\begin{document}
\maketitle
\vspace{-0.5cm}

\begin{abstract}
\noindent
Ce rapport présente un benchmark de performance de VASP 6.5.0 sur le cluster de calcul local,
réalisé en amont de la campagne de relaxation des polymorphes Al$_4$C$_3$. L'objectif est de
déterminer, pour des calculs de type SCF (point unique), la combinaison de paramètres de
parallélisation (\texttt{NPAR}, \texttt{KPAR}) et le nombre de tâches MPI qui minimisent le temps
CPU, sur trois systèmes de taille et de densité de points \textbf{k} différentes. Le meilleur
compromis obtenu est \texttt{NPAR=1} avec \texttt{KPAR} égal à 4 ou 8 selon le système, apportant
un gain allant de $\times1.5$ à $\times7.5$ par rapport à la configuration par défaut
(\texttt{NPAR=KPAR=1}), tandis que l'augmentation seule du nombre de tâches MPI sans découpage en
points \textbf{k} sature rapidement.
\end{abstract}

\section{Objectifs}

Avant de lancer les relaxations ioniques des polymorphes Al$_4$C$_3$ (calculs coûteux, jusqu'à
28 atomes, cellules à 0 et 50~GPa), un benchmark de performance a été mené pour identifier une
configuration de parallélisation efficace sur le cluster SLURM local. Deux questions sont posées :

\begin{enumerate}[label=(\alph*)]
  \item \textbf{Partie A} — à nombre de tâches MPI fixe (\texttt{ntasks=40}), quelle combinaison
  \texttt{NPAR}/\texttt{KPAR} minimise le temps de calcul, et cette combinaison optimale
  dépend-elle du système traité ?
  \item \textbf{Partie B} — à parallélisation par bandes/points \textbf{k} désactivée
  (\texttt{NPAR=KPAR=1}), comment le temps de calcul évolue-t-il avec le nombre de tâches MPI
  seul (4 à 40) ?
\end{enumerate}

\section{Configuration matérielle et logicielle}

\subsection{Cluster}

Les calculs ont été soumis via SLURM (partition \texttt{defq}) sur des nœuds mono-nœud
(\texttt{--nodes=1}), avec deux familles de nœuds au matériel distinct présentes sur le cluster :

\begin{table}[h]
\centering
\begin{tabular}{lccc}
\toprule
Famille & Nœuds & Configuration & CPU logiques \\
\midrule
\texttt{48c\_2thread} & node01--node08 & 2 sockets $\times$ 12 c\oe urs $\times$ 2 threads (SMT) & 48 \\
\texttt{40c\_1thread} & node09--node16 & 2 sockets $\times$ 20 c\oe urs $\times$ 1 thread (SMT off) & 40 \\
\bottomrule
\end{tabular}
\caption{Familles de nœuds de calcul utilisées dans le benchmark.}
\end{table}

Le nœud d'exécution n'était pas fixé explicitement dans les scripts de soumission ; l'ordonnanceur
SLURM a donc réparti les runs entre les deux familles (colonne \texttt{node\_family} du fichier de
résultats). Ce point est repris comme limite méthodologique en section~\ref{sec:limites}.

\subsection{Code et paramètres communs}

VASP 6.5.0 (\texttt{vasp\_std}, mode PAW), MPI Intel (\texttt{impi/2021.13}), un thread OpenMP par
tâche (\texttt{OMP\_NUM\_THREADS=1}). Tous les calculs sont des points SCF simples
(\texttt{IBRION=-1}, \texttt{NSW=0}, pas de relaxation ionique), avec les paramètres communs
suivants~:

\begin{center}
\begin{tabular}{ll}
\toprule
Paramètre & Valeur \\
\midrule
\texttt{PREC} & Accurate \\
\texttt{ENCUT} & 600~eV \\
\texttt{EDIFF} & $10^{-5}$ \\
\texttt{ISMEAR / SIGMA} & 0 / 0.05 \\
\texttt{LWAVE / LCHARG} & .FALSE. / .FALSE. \\
\bottomrule
\end{tabular}
\end{center}

\subsection{Systèmes de test}

Trois systèmes de taille croissante et de maillages \textbf{k} très différents (à
\texttt{KSPACING} fixé par système) ont été utilisés, afin de vérifier si la configuration
optimale dépend du nombre de points \textbf{k} irréductibles :

\begin{table}[h]
\centering
\begin{tabular}{lcccc}
\toprule
Système & Atomes & Formule & \texttt{KSPACING} (\AA$^{-1}$) & Points \textbf{k} irréductibles \\
\midrule
\texttt{Ta4C3\_7at}   & 7  & Ta$_4$C$_3$   & 0.03 & 189 \\
\texttt{Ccubic\_16at} & 16 & C$_{16}$      & 0.03 & 105 \\
\texttt{As4S3\_28at}  & 28 & As$_{16}$S$_{12}$ & 0.20 & 27 \\
\bottomrule
\end{tabular}
\caption{Systèmes de test, densité de maillage \textbf{k} et nombre de points \textbf{k}
irréductibles résultant (grille de Monkhorst--Pack centrée $\Gamma$).}
\end{table}

Le \texttt{KSPACING} de \texttt{As4S3\_28at} est volontairement plus lâche (0.20 contre 0.03) et a
été gardé identique sur tous ses runs (Parties A et B), de sorte que les 10 mesures sur ce système
restent directement comparables entre elles.

\section{Méthodologie}

Les runs ont été générés automatiquement (\texttt{gen\_runs.py}) : POSCAR/POTCAR communs par
système, INCAR et script SLURM générés par combinaison, limite de temps de 30~min par job.
22~jobs ont été soumis (IDs SLURM 58341--58362) et suivis jusqu'à completion via une boucle de
polling (\texttt{wait\_loop.sh}, intervalle 60~s).

\begin{itemize}
  \item \textbf{Partie A} (18 runs) : les 3 systèmes $\times$ 6 combinaisons
  \texttt{(NPAR, KPAR)} $\in \{(1,1),(4,2),(2,4),(1,4),(8,1),(1,8)\}$, à \texttt{ntasks=40} fixe.
  \item \textbf{Partie B} (4 runs) : le système \texttt{As4S3\_28at} seul, à
  \texttt{NPAR=KPAR=1} fixe, pour \texttt{ntasks} $\in \{4,8,16,24\}$ — complété par le point
  \texttt{ntasks=40, NPAR=KPAR=1} déjà mesuré en Partie A.
\end{itemize}

Le temps CPU (\texttt{cpu\_time\_s}) et le temps écoulé (\texttt{elapsed\_time\_s}) sont extraits de
la sortie \texttt{time -p} du wrapper \texttt{srun}. Chaque configuration n'a été mesurée
\textbf{qu'une seule fois} (pas de répétition ni d'écart-type~; voir limites,
section~\ref{sec:limites}).

\section{Résultats}

\subsection{Partie A --- sweep NPAR/KPAR à \texttt{ntasks=40}}

\begin{table}[h]
\centering
\begin{tabular}{lccS[table-format=3.1]S[table-format=1.2]}
\toprule
Système & NPAR & KPAR & {Temps CPU (s)} & {Accélération} \\
\midrule
Ta4C3\_7at   & 1 & 1 & 437.9 & 1.00 \\
Ta4C3\_7at   & 4 & 2 & 268.4 & 1.63 \\
Ta4C3\_7at   & 2 & 4 & 267.6 & 1.64 \\
Ta4C3\_7at   & 1 & 4 & \bfseries 232.0 & \bfseries 1.89 \\
Ta4C3\_7at   & 8 & 1 & 369.9 & 1.18 \\
Ta4C3\_7at   & 1 & 8 & 277.5 & 1.58 \\
\midrule
Ccubic\_16at & 1 & 1 & 265.7 & 1.00 \\
Ccubic\_16at & 4 & 2 & 46.1  & 5.77 \\
Ccubic\_16at & 2 & 4 & 39.2  & 6.77 \\
Ccubic\_16at & 1 & 4 & 44.2  & 6.01 \\
Ccubic\_16at & 8 & 1 & 69.7  & 3.81 \\
Ccubic\_16at & 1 & 8 & \bfseries 35.5  & \bfseries 7.49 \\
\midrule
As4S3\_28at  & 1 & 1 & 372.5 & 1.00 \\
As4S3\_28at  & 4 & 2 & 339.7 & 1.10 \\
As4S3\_28at  & 2 & 4 & 313.2 & 1.19 \\
As4S3\_28at  & 1 & 4 & \bfseries 240.6 & \bfseries 1.55 \\
As4S3\_28at  & 8 & 1 & 375.6 & 0.99 \\
As4S3\_28at  & 1 & 8 & 330.8 & 1.13 \\
\bottomrule
\end{tabular}
\caption{Temps CPU et accélération par rapport à \texttt{NPAR=KPAR=1}, à \texttt{ntasks=40} fixe.
La meilleure combinaison par système est en gras.}
\label{tab:partA}
\end{table}

\begin{figure}[h]
\centering
\begin{tikzpicture}
\begin{axis}[
  width=\textwidth, height=7cm,
  ybar, bar width=5pt,
  ymode=log,
  ylabel={Temps CPU (s), échelle log},
  symbolic x coords={1-1,4-2,2-4,1-4,8-1,1-8},
  xtick=data,
  xlabel={NPAR-KPAR},
  legend pos=outer north east,
  legend style={font=\small},
  enlarge x limits=0.12,
  nodes near coords style={font=\tiny},
]
\addplot coordinates {(1-1,437.929) (4-2,268.370) (2-4,267.591) (1-4,231.992) (8-1,369.869) (1-8,277.543)};
\addplot coordinates {(1-1,265.650) (4-2,46.068) (2-4,39.241) (1-4,44.227) (8-1,69.666) (1-8,35.479)};
\addplot coordinates {(1-1,372.519) (4-2,339.687) (2-4,313.182) (1-4,240.605) (8-1,375.632) (1-8,330.788)};
\legend{Ta4C3\_7at (189 k-pts), Ccubic\_16at (105 k-pts), As4S3\_28at (27 k-pts)}
\end{axis}
\end{tikzpicture}
\caption{Temps CPU (échelle logarithmique) pour les 6 combinaisons NPAR/KPAR testées, à
\texttt{ntasks=40}, pour les 3 systèmes. Le gain lié à \texttt{KPAR>1} est d'autant plus marqué
que le système a beaucoup de points \textbf{k} irréductibles.}
\label{fig:partA}
\end{figure}

\subsection{Partie B --- scaling en nombre de tâches MPI (As4S3\_28at, NPAR=KPAR=1)}

\begin{table}[h]
\centering
\begin{tabular}{cSSlc}
\toprule
{ntasks} & {Temps CPU (s)} & {Accél. vs 4} & Nœud & Famille \\
\midrule
4  & 851.654 & 1.00 & node15 & 40c\_1thread \\
8  & 536.230 & 1.59 & node16 & 40c\_1thread \\
16 & 539.682 & 1.58 & node04 & 48c\_2thread \\
24 & 408.179 & 2.09 & node04 & 48c\_2thread \\
40 & 372.519 & 2.29 & node04 & 48c\_2thread \\
\bottomrule
\end{tabular}
\caption{Temps CPU en fonction du nombre de tâches MPI, à \texttt{NPAR=KPAR=1} (pas de
parallélisation par points \textbf{k} ni par bandes). Le changement de famille de nœud entre
\texttt{ntasks=8} et \texttt{ntasks=16} est indiqué (voir section~\ref{sec:limites}).}
\label{tab:partB}
\end{table}

\begin{figure}[h]
\centering
\begin{tikzpicture}
\begin{axis}[
  width=\textwidth, height=7cm,
  xlabel={Nombre de tâches MPI (ntasks)},
  ylabel={Temps CPU (s)},
  xmode=log, log basis x=2,
  xtick={4,8,16,24,40},
  xticklabels={4,8,16,24,40},
  legend pos=north east,
  legend style={font=\small},
  grid=major,
]
\addplot[mark=*, thick, color=blue!70!black]
  coordinates {(4,851.654) (8,536.230) (16,539.682) (24,408.179) (40,372.519)};
\addlegendentry{Mesuré (NPAR=KPAR=1)}
\addplot[mark=none, dashed, color=gray]
  coordinates {(4,851.654) (40,85.1654)};
\addlegendentry{Scaling linéaire idéal}
\end{axis}
\end{tikzpicture}
\caption{Temps CPU vs nombre de tâches MPI pour \texttt{As4S3\_28at}, comparé au scaling linéaire
idéal extrapolé depuis \texttt{ntasks=4}. Le décrochage est net dès \texttt{ntasks=8} : sans
découpage en points \textbf{k} (\texttt{KPAR=1}), le système (27 points \textbf{k} irréductibles)
sature bien avant \texttt{ntasks=40}.}
\label{fig:partB}
\end{figure}

\section{Discussion}

\subsection{Effet de NPAR/KPAR}

Sur les trois systèmes, \texttt{KPAR=4} ou \texttt{KPAR=8} avec \texttt{NPAR=1} domine
systématiquement les combinaisons à \texttt{NPAR>1} (tableau~\ref{tab:partA}, figure~\ref{fig:partA}) :

\begin{itemize}
  \item \textbf{Ccubic\_16at} (105 points \textbf{k}) : gain maximal $\times7.49$ avec
  \texttt{KPAR=8} — le système a largement assez de points \textbf{k} pour remplir 8 groupes.
  \item \textbf{Ta4C3\_7at} (189 points \textbf{k}) : gain maximal $\times1.89$ seulement avec
  \texttt{KPAR=4}, malgré un nombre de points \textbf{k} très élevé — ce système, plus petit
  (7 atomes), est probablement limité par un autre facteur (base d'ondes planes réduite, overhead
  de communication) plutôt que par le découpage \textbf{k}.
  \item \textbf{As4S3\_28at} (27 points \textbf{k}) : gain maximal $\times1.55$ avec
  \texttt{KPAR=4} ; \texttt{KPAR=8} est ici \emph{moins} bon (1.13$\times$) que \texttt{KPAR=4}
  (1.55$\times$), et \texttt{NPAR=8/KPAR=1} n'apporte aucun gain (0.99$\times$). Avec seulement
  27 points \textbf{k}, répartis sur 8 groupes MPI, chaque groupe ne reçoit que 3--4 points
  \textbf{k} et se retrouve avec peu de tâches par groupe (5 rangs/groupe) : l'overhead de
  communication au sein de chaque groupe l'emporte sur le gain de parallélisation en \textbf{k}.
\end{itemize}

\textbf{Aucune combinaison à \texttt{NPAR>1} testée (4/2, 2/4, 8/1) ne bat la meilleure
combinaison à \texttt{NPAR=1}} sur un seul des trois systèmes. La parallélisation par bandes
(\texttt{NPAR}) semble ici globalement moins efficace que la parallélisation par points
\textbf{k} (\texttt{KPAR}) pour ces tailles de système et ce nombre de tâches.

\subsection{Effet du nombre de tâches MPI seul}

Sans découpage \textbf{k} (\texttt{NPAR=KPAR=1}, tableau~\ref{tab:partB},
figure~\ref{fig:partB}), le passage de 4 à 40 tâches ne procure qu'un facteur $\times2.29$, très
loin du $\times10$ idéal. L'essentiel du décrochage intervient déjà entre \texttt{ntasks=4} et
\texttt{ntasks=8} (efficacité de $\sim$79~\% en doublant les tâches), puis le gain devient marginal
au-delà de 16 tâches. Ceci confirme, en creux, l'intérêt de la Partie A : passer de 40 tâches
« brutes » à 40 tâches réparties en \texttt{KPAR=4} groupes de points \textbf{k} rapporte plus
que d'augmenter le nombre de tâches MPI sans structurer le parallélisme.

\section{Recommandations pratiques}

Pour les calculs de relaxation Al$_4$C$_3$ (28 atomes, \texttt{ENCUT=600}, cellules à 0 et
50~GPa, \texttt{KSPACING=0.03}, donc un nombre de points \textbf{k} irréductibles probablement
élevé comme \texttt{Ta4C3\_7at}/\texttt{Ccubic\_16at} plutôt que comme
\texttt{As4S3\_28at}) :

\begin{itemize}
  \item Utiliser \texttt{NPAR=1} par défaut, et régler \texttt{KPAR} en fonction du nombre de
  points \textbf{k} irréductibles réel du système (\texttt{KPAR=4} est un choix sûr et
  quasi-optimal dans les trois cas testés ; \texttt{KPAR=8} n'aide que si le système a
  suffisamment de points \textbf{k} pour remplir 8 groupes sans sous-alimenter chaque groupe).
  \item Ne pas dépasser \texttt{ntasks=24--40} sans \texttt{KPAR>1} : au-delà, le gain marginal
  est faible pour des systèmes de taille comparable (16--28 atomes).
  \item Vérifier le nombre de points \textbf{k} irréductibles (ligne \texttt{"irreducible
  k-points"} de l'\texttt{OUTCAR}) avant de figer \texttt{KPAR} pour un nouveau polymorphe, plutôt
  que de réutiliser aveuglément la config optimale d'un autre système.
\end{itemize}

C'est la configuration \texttt{NPAR=1 / KPAR=4} qui a été retenue dans les scripts de relaxation
Al$_4$C$_3$ (\texttt{vasp6\_relax.sh}, cf.\ \texttt{INCAR.relax}) suite à ce benchmark.

\section{Limites de l'étude}
\label{sec:limites}

\begin{itemize}
  \item \textbf{Une seule mesure par configuration} : aucune répétition n'a été effectuée, donc
  aucune barre d'incertitude n'est disponible. La variabilité liée à la charge du nœud ou au
  réseau n'est pas quantifiée.
  \item \textbf{Hétérogénéité matérielle non contrôlée en Partie B} : les points
  \texttt{ntasks=4} et \texttt{ntasks=8} ont tourné sur des nœuds sans SMT
  (\texttt{40c\_1thread}), tandis que \texttt{ntasks=16, 24, 40} ont tourné sur un nœud avec SMT
  actif (\texttt{48c\_2thread}). La comparaison brute entre ces deux groupes de points n'isole
  donc pas purement l'effet du nombre de tâches.
  \item \textbf{Systèmes non représentatifs à 100\%} des polymorphes Al$_4$C$_3$ : les trois
  systèmes de test diffèrent en composition chimique et en \texttt{KSPACING}. Les conclusions sur
  la meilleure combinaison \texttt{NPAR}/\texttt{KPAR} sont extrapolées par analogie sur le
  nombre de points \textbf{k}, pas vérifiées directement sur Al$_4$C$_3$.
  \item Les temps rapportés sont des temps CPU/écoulés pour un unique point SCF
  (\texttt{IBRION=-1}), non pour une relaxation ionique complète ; l'INCAR de production diffère
  légèrement (\texttt{IBRION=2}, \texttt{ISIF=8}, \texttt{EDIFFG=-0.01}).
\end{itemize}

\section{Conclusion}

Le découpage par points \textbf{k} (\texttt{KPAR}) est, sur les trois systèmes testés, nettement
plus efficace que le découpage par bandes (\texttt{NPAR}) ou que la simple augmentation du nombre
de tâches MPI. La configuration \texttt{NPAR=1 / KPAR=4} constitue un choix robuste par défaut
(gain $\times1.5$ à $\times6$ selon le système), avec un gain supplémentaire possible via
\texttt{KPAR=8} uniquement pour les systèmes disposant d'un grand nombre de points \textbf{k}
irréductibles. Ce réglage a été adopté pour la campagne de relaxation des polymorphes Al$_4$C$_3$.

\appendix
\section{Détails de soumission}

22 jobs soumis sur SLURM (partition \texttt{defq}), IDs 58341 à 58362, générés et suivis via
\texttt{benchmarks/scf\_speed\_test/gen\_runs.py} et \texttt{wait\_loop.sh}. Données brutes :
\texttt{benchmarks/scf\_speed\_test/results\_cpu\_time.csv}.

\end{document}
